\chapter{Sai, quê và cảm giác thua rất thật}
\markboth{Sai, quê và cảm giác thua rất thật}{Sai, quê và cảm giác thua rất thật}

\section{Em ngại phát biểu vì sai là bị cười}
\begin{truyen}{Em ngại phát biểu vì sai là bị cười}{Classroom Talk}
\chuhoa{M}{ột bạn nói nhỏ nhưng thật: ``Em không phải không biết gì. Em chỉ sợ nói sai xong bị cười, rồi quê tới hết buổi.''}

Thầy hỏi: ``Vậy em muốn quê một lần để biết đường sửa, hay dốt âm thầm lâu dài?''

Bạn ấy cười méo: ``Thầy nói câu nào cũng đâm trúng chỗ em né.''

Thầy đáp: ``Vì chỗ em né lại chính là chỗ em cần bước qua.''

\textit{(Fear of embarrassment keeps many students quiet, but silence can preserve weakness far longer than one awkward moment.)}
\end{truyen}

\begin{baihoc}
Nỗi quê ngắn hạn thường rẻ hơn cái giá của việc im lặng quá lâu.
\textit{(Short-term embarrassment often costs less than long-term silence.)}
\end{baihoc}

\begin{ghinhoanh}
\textbf{Short English to remember:}

Embarrassment passes.

Untouched weakness stays.
\end{ghinhoanh}

\begin{tuhocnhanh}
1. Em đang né học vì sợ sai hay sợ bị cười? \textit{(Are you avoiding learning because of error or embarrassment?)}

2. Em có thể tập phát biểu ở môi trường an toàn hơn trước không? \textit{(Can you practice speaking in a safer setting first?)}
\end{tuhocnhanh}
\ngancach

\section{Thua đứa mình từng giúp còn đau hơn}
\begin{truyen}{Thua đứa mình từng giúp còn đau hơn}{Classroom Talk}
\chuhoa{C}{ó bạn chua chát: ``Thầy ơi, thua người giỏi sẵn em chịu. Nhưng thua đứa mình từng chỉ bài mới cay.''}

Thầy hỏi: ``Em cay vì nó giỏi lên, hay cay vì em tưởng nó phải đứng dưới mình mãi?''

Bạn ấy im, rồi thú nhận: ``Chắc cả hai.''

Thầy nói: ``Vậy em đau không chỉ vì kết quả. Em đau vì cái tôi bị trầy.''

\textit{(Losing to someone we once helped can expose more pride than injustice.)}
\end{truyen}

\begin{baihoc}
Người khác tiến bộ không phải là phản bội mình. Đôi khi chỉ là mình đã đứng yên quá lâu.
\textit{(Another person improving is not a betrayal. Sometimes it only shows that we stood still too long.)}
\end{baihoc}

\begin{ghinhoanh}
\textbf{Short English to remember:}

Their growth is not your insult.

Pride often hurts before reality does.
\end{ghinhoanh}

\begin{tuhocnhanh}
1. Em đang buồn vì thua, hay vì cái tôi bị chạm? \textit{(Are you upset by the loss, or by wounded pride?)}

2. Em cần học gì từ người đã vượt lên? \textit{(What can you learn from the person who surpassed you?)}
\end{tuhocnhanh}
\ngancach

\section{Em sửa mấy lần vẫn sai nên muốn bỏ}
\begin{truyen}{Em sửa mấy lần vẫn sai nên muốn bỏ}{Classroom Talk}
\chuhoa{M}{ột bạn đập bút xuống bàn: ``Em sửa đi sửa lại vẫn trật. Cảm giác như đầu em không hấp thụ nổi nữa. Muốn bỏ luôn cho xong.''}

Thầy nói: ``Bỏ lúc này sẽ cho em cảm giác nhẹ ngay. Nhưng cái sai đó không tự biến mất. Nó chỉ chờ dịp khác quay lại đập em mạnh hơn.''

Bạn hỏi: ``Vậy cứ lì tiếp hả thầy?''

Thầy đáp: ``Không phải lì mù. Là đổi cách sửa, đổi người hỏi, đổi góc nhìn. Đừng cứ dùng một kiểu rồi kết luận mình hết cứu.''

\textit{(Repeated failure does not always mean incapacity. Sometimes it means the current method has reached its limit.)}
\end{truyen}

\begin{baihoc}
Muốn bỏ thường đến ngay trước lúc cần đổi cách, không phải lúc cần tự kết thúc mình.
\textit{(The urge to quit often appears right when a change of method is needed, not a final surrender.)}
\end{baihoc}

\begin{ghinhoanh}
\textbf{Short English to remember:}

Repeating the same failed method is not resilience.

Change the method before judging yourself.
\end{ghinhoanh}

\begin{tuhocnhanh}
1. Em đã thử đổi cách học ở chỗ sai đó chưa? \textit{(Have you changed your method for that weak point?)}

2. Ai có thể giúp em nhìn ra lỗi theo cách khác? \textit{(Who can help you see the mistake differently?)}
\end{tuhocnhanh}
\ngancach

\section{Có lúc em thấy mình cố cũng chỉ vậy}
\begin{truyen}{Có lúc em thấy mình cố cũng chỉ vậy}{Classroom Talk}
\chuhoa{C}{ó bạn nói câu nghe rất mỏi: ``Thầy ơi, có lúc em thấy mình cố rồi mà đời vẫn không đổi bao nhiêu. Nên em tự hỏi cố nữa để làm gì.''}

Thầy đáp chậm rãi: ``Có những giai đoạn cố gắng chưa ra quả ngay. Nhưng nó đang âm thầm dời cái nền của em. Vấn đề là mình có đủ kiên nhẫn để chịu đoạn chưa thấy gì không.''

Bạn ấy nhìn xuống: ``Đoạn đó khó chịu thật.''

Thầy nói: ``Ừ. Trưởng thành là vẫn làm việc đúng trong lúc chưa được thưởng.''

\textit{(Some seasons of effort look invisible. Yet deep change often begins before visible reward appears.)}
\end{truyen}

\begin{baihoc}
Cái khó nhất không phải cố một bữa. Là cố qua giai đoạn chưa được công nhận.
\textit{(The hardest part is not trying once. It is trying through the phase where no one notices.)}
\end{baihoc}

\begin{ghinhoanh}
\textbf{Short English to remember:}

Invisible effort still matters.

Growth often starts below the surface.
\end{ghinhoanh}

\begin{tuhocnhanh}
1. Em có đang bỏ một việc đúng chỉ vì chưa thấy kết quả sớm? \textit{(Are you dropping the right work just because results are slow?)}

2. Em cần thêm bằng chứng nào để tin vào quá trình của mình? \textit{(What evidence would help you trust your process more?)}
\end{tuhocnhanh}
\ngancach